\clearpage
\subsection{G7：DI 原生请求当前可达状态}
\begin{center}
\begin{tikzpicture}
\node[dstate] (new) at (0,0) {request 参数校验\\创建 Operation};
\node[dstate] (pending) at (5,0) {Pending\\固定绝对 deadline};
\node[dstate] (prep) at (10,0) {PreparingInput\\PreparedModel 输入投影};
\node[dstate] (core) at (10,-3) {Core request\\ACK / Selection /\\Response};
\node[dstate] (fail) at (5,-6) {Failed\\兼容构造无 preparation/contract};
\node[dstate] (cancel) at (5,-3) {Cancelled\\cancel 竞争终态};
\node[dstate] (timeout) at (0,-3) {Failed\\请求 deadline 到期};
\draw[down] (new)--(pending);
\draw[down] (pending)--(prep);
\draw[down] (prep)--node[dlabel,right]{prepared runtime}(core);
\draw[down] (pending)--node[dlabel,left]{无 runtime/contract}(fail);
\draw[down] (pending)--node[dlabel,right]{尚未终态}(cancel);
\draw[down] (pending)--node[dlabel,above,sloped]{超时先发生}(timeout);
\end{tikzpicture}
\end{center}
\diagramnote{图混合展示内部 phase（PreparingInput）与对外 status（Pending/Failed/Cancelled），
它们不是同一个枚举。PreparedModel 路径从输入投影进入 Core request，并由 C++ 过程矩阵验证
ACK/Selection/Response；没有 preparation/request contract 的兼容构造保留
\code{NATIVE_REQUEST_PIPELINE_NOT_READY}。}
\callout{终态与等待}{cancel、deadline 和 dispatch 失败竞争一次终态写入；已有终态不被后续路径覆盖。
result(waitTimeout) 的局部超时只结束这次等待，不自动取消 operation。构造或提交前的非法参数异常未画成异步终态。}
\diagramnote{源码：\code{PreparedModel.cpp} 与 \code{NativeInferenceClient.cpp} 的 request、dispatchOperation、
markTerminal、cancelOperation 和 result。声明中存在 Succeeded 等状态不代表默认路径已可达。
Python facade 只转换参数、异常、GIL 与 asyncio，不另建请求状态机。}
